\section{Morphisms of root data}\label{sga3:XXI:6}
\renewcommand{\thetheorem}{\thesubsection.\arabic{theorem}}

\subsection{Definition}\label{sga3:XXI:6.1}

Let
\[
 \cR=(M,M^\ast,R,R^\ast),
 \qquad
 \cR'=(M',M'^\ast,R',R'^\ast)
\]
be two root data. Let \(f:M'\to M\) be a linear map and
\({}^tf:M^\ast\to M'^\ast\) its transpose.

\begin{definition}\label{sga3:XXI:6.1.1}
The map \(f\) is called a \emph{morphism} from \(\cR'\) to \(\cR\),
and is denoted
\[
 f:\cR'\longrightarrow\cR,
\]
if \(f\) induces a bijection from \(R'\) onto \(R\), and
\({}^tf\) induces a bijection from \(R^\ast\) onto \(R'^\ast\).
\end{definition}

Then \({}^tf\) is a morphism of the dual root data:
\[
 {}^tf:\cR^\ast\longrightarrow\cR'^\ast.
\]
If \(f:\cR'\to\cR\) is a morphism and
\(\alpha=f(\alpha')\) for \(\alpha'\in R'\), then
\[
 \alpha'^\ast={}^tf(\alpha^\ast).
\]
Indeed, if \(p\) and \(p'\) denote the maps of
\ref{sga3:XXI:1.2.1} for \(\cR\) and \(\cR'\), respectively, then
\[
 p'={}^tf\circ p\circ f,
\]
which gives the assertion at once. The reader may verify the
statements below, almost all of which are immediate.

\begin{proposition}\label{sga3:XXI:6.1.2}
Let \(f:\cR'\to\cR\) be a morphism of root data. If
\(\alpha'\in R'\) and \(\alpha=f(\alpha')\), then
\[
 \alpha'^\ast={}^tf(\alpha^\ast).
\]
Moreover, \(f\) induces isomorphisms
\[
 R'\xrightarrow{\sim}R,\qquad
 \Gamma_0(R')\xrightarrow{\sim}\Gamma_0(R),\qquad
 \VR(R')\xrightarrow{\sim}\VR(R),
\]
and \({}^tf\) induces isomorphisms
\[
 R^\ast\xrightarrow{\sim}R'^\ast,\qquad
 \Gamma_0(R^\ast)\xrightarrow{\sim}\Gamma_0(R'^\ast),\qquad
 \VR(R^\ast)\xrightarrow{\sim}\VR(R'^\ast).
\]
The last is the transpose of the corresponding morphism induced by
\(f\). The map
\[
 s_{\alpha'}\longmapsto s_{f(\alpha')}
\]
extends to an isomorphism
\[
 W(\cR')\xrightarrow{\sim}W(\cR)
\]
compatible with the actions of these two groups on the sets of
\ref{sga3:XXI:1.1.13}.
\end{proposition}

\begin{proposition}\label{sga3:XXI:6.1.3}
The maps
\[
 \Delta'\longmapsto f(\Delta'),\qquad
 R'_+\longmapsto f(R'_+),\qquad
 C'\longmapsto({}^tf\otimes\mathbb R)^{-1}(C')
\]
define bijective correspondences between the systems of simple roots,
the systems of positive roots, and the Weyl chambers for \(\cR'\) and
\(\cR\). These correspondences are compatible with the actions of the
Weyl groups and with the correspondences
\[
 \mathscr{S}(R_+)\longleftrightarrow R_+
 \longleftrightarrow\mathscr{C}(R_+).
\]
\end{proposition}

\begin{lemma}\label{sga3:XXI:6.1.4}
Morphisms compose. A morphism \(f:\cR'\to\cR\) is an isomorphism if
and only if the linear map \(f:M'\to M\) is bijective.
\end{lemma}
